
\section{Introduction}
\label{sec:intro}

% Overall
Rolling element bearings let parts rotate with minimal friction in almost every kind of rotating machine: motors, pumps, wind turbines, machine tools, trains. What makes this possible is a thin film of oil or grease, that keeps the rolling balls from touching their metal tracks. As long as that film is intact, the bearing runs cool with low friction and wears slowly. Once it breaks down -- because there is too little lubricant, the lubricant has degraded, or it has become contaminated -- metal starts touching metal, friction rises, heat builds up, and wear accelerates sharply, ultimately causing premature mechanical failure. Lubrication problems are estimated to cause 30--80\% of all rolling element bearing failures, more than mechanical overload or installation errors combined \cite{jakobsen_detecting_2023,dai_situ_2024,xu_vibration-based_2024,nassef_impact_2022,isalm_oil_2021}.

% Consequences of failures
The cost of this is substantial. A single bearing failure can shut down an entire production line, and the resulting downtime typically costs far more than the bearing itself. The problem is especially acute where access is difficult: at an offshore wind farm, replacing one failed bearing can require a vessel and a crane crew, and a single unplanned trip can cost more than the part it replaces several times over \cite{wakiru_review_2019}. Reducing lubrication-related failures is therefore a large available opportunity for improving machine uptime and cutting maintenance spend across industry, as well as cutting energy needed to overcome friction.

% Existing methods are insufficient
Despite this, most industrial lubrication schedules are still time-based: lubricant is topped up or replaced at a fixed interval, regardless of whether the bearing actually needs it at that moment. Re-lubricating too early wastes lubricant, raises operating temperature, and increases friction. Re-lubricating too late lets the film thin out and wear begin. The correct interval depends on bearing geometry, speed, load, temperature, and the lubricant's own condition -- a moving target that no fixed schedule can track \cite{wakiru_review_2019}.

\textit{Condition-based relubrication} (CBR) replaces the schedule with a measurement: lubricate when a sensor signal shows it is actually needed, not before. CBR depends on \textit{lubrication condition monitoring} (LCM): continuously estimating whether the lubricant film is thick enough, how degraded the lubricant has become, and whether it is contaminated. Reliable LCM would let maintenance teams act on the bearing's actual condition rather than a guess.

LCM has received far less research attention than the related and more mature problem of detecting mechanical bearing faults from vibration, which is now standard practice in industry. This is not because the necessary sensors are missing: a range of sensing technologies capable of picking up lubrication-related signals -- acoustic emission, ultrasound, vibration, electrical capacitance, and others -- have already been demonstrated in laboratory studies and ultrasound is used with simple relative threshold methods in industrial applications. The bottleneck is what happens after sensing: turning a raw sensor signal into a reliable estimate of lubrication condition that still works once the bearing's speed, load, or temperature changes. This signal processing step is fragmented across sensing technologies and rarely tested outside the lab. This gap is the subject of this review.

Throughout this review, \textit{online} monitoring means continuous sensing on a bearing that keeps running, with no need to stop it or extract lubricant; \textit{in-line} monitoring means sensors built into the lubricant circulation system, which also requires no stoppage. Both are in scope. Methods that need a brief interruption (\textit{at-line}) or a lab sample (\textit{offline}) are not.

\subsection{Existing Reviews and Their Limitations}
\label{subsec:existing_reviews}

A number of reviews on LCM exists: Table~\ref{tab:existing_reviews} lists the reviews most closely related to this one and what each covers. \todo{Change table form to prose. Check Wakiru review from MSSP for style.}

\begin{table}[H]
\small
\centering
\renewcommand{\arraystretch}{1.3}
\begin{tabularx}{\linewidth}{ >{\raggedright\arraybackslash}p{3.2cm} >{\raggedright\arraybackslash}p{4cm} >{\raggedright\arraybackslash}X }
\hline
\textbf{Review} & \textbf{Focus} & \textbf{Limitation relative to this review} \\
\hline
Wakiru et al.\ \cite{wakiru_review_2019} & LCM for maintenance decision support & Centred on invasive lubricant sampling and analysis, not online sensing \\
Shang et al.\ \cite{shang_comprehensive_2024} & Film thickness measurement & Active/invasive methods only (optical, electrical, ultrasonic) \\
Wang et al.\ \cite{wang_advanced_2024} & Bearing condition monitoring, incl.\ film thickness and debris & Mostly active/invasive; passive methods mentioned only in passing \\
Hansen et al.\ \cite{hansen_comparison_2022} & Ultrasonic reflectometry calibration & Single active method only \\
Cen et al.\ \cite{cen_ehl_2021}; Becker-Dombrowsky et al.\ \cite{becker-dombrowsky_electrical_2024} & Electrical capacitance / impedance & Single active method only \\
Dai et al.\ \cite{dai_situ_2024} & In-situ film thickness by capacitance & Single active method only \\
Zhu et al.\ \cite{zhu_lubricating_2017,zhu_survey_2013} & Online lubricant sensors, oil monitoring & Active/invasive methods only \\
Sun et al.\ \cite{sun_online_2021} & Wear debris monitoring & Active methods only \\
Dou et al.\ \cite{dou_review_2022} & Ultrasonic film thickness & Single active method only \\
\hline
\end{tabularx}
\caption{Reviews related to lubrication condition monitoring and how their scope differs from this review.}
\label{tab:existing_reviews}
\end{table}

Three limitations are shared by all of them. First, none organises the field by signal processing method -- signal processing is treated as a detail beneath each sensing technology, not as a subject in its own right. Second, almost all focus on active or invasive sensing, leaving passive, non-invasive methods underrepresented even though these offer the easiest path to industrial retrofitting without halting machine operation. Third, none identifies signal processing -- rather than sensor hardware -- as the main reason LCM has not moved from the lab into industrial deployment.

\subsection{Contribution and Scope of This Review}
\label{subsec:contribution}

This review makes four contributions:
\begin{enumerate}
    \item \textbf{A systematic taxonomy of signal processing methods for LCM}, organised by method family and applied across all sensing technologies -- the first such taxonomy in the LCM literature.
    \item \textbf{A structured comparison} of sensing and signal processing combinations across four dimensions: sensitivity to different lubrication faults, sensitivity to changing operating conditions, practical deployment requirements, and technology maturity.
    \item \textbf{A clear list of signal-processing-specific research gaps}, each with a concrete suggested research direction.
    \item \textbf{Coverage of all sensing technologies} -- active and passive, invasive and non-invasive -- with extra depth on the passive, non-invasive methods most relevant to retrofitting existing industrial bearings.
\end{enumerate}
\todo{Update in the end if contributions are changed}

This review covers online and in-line LCM literature for rolling element bearings (primarily ball bearings), restricted to studies that describe their signal processing method in enough detail to assess its assumptions and limitations. Studies on bearing fault detection with no lubrication-specific focus, and studies on journal bearings or other non-bearing tribological contacts, are excluded. The literature search covers publications up to July 2026.

\subsection{Paper Organisation}
\label{subsec:organisation}

Section~\ref{sec:search} describes the search and inclusion criteria. Section~\ref{sec:lcm_pipeline} defines the online LCM pipeline and establishes the scope of this review within it. Section~\ref{sec:lubrication_physics} introduces the relevant lubrication physics and degradation mechanisms. Section~\ref{sec:signal_origins} explains how lubrication condition shows up physically in sensor signals. Section~\ref{sec:sensing} surveys sensing technologies, including the active/passive and invasive/non-invasive classification used throughout. Section~\ref{sec:sp} -- the main contribution -- reviews signal processing methods by family. Section~\ref{sec:comparative} compares sensing and signal processing combinations. Section~\ref{sec:gaps} sets out research gaps and future directions, and Section~\ref{sec:conclusion} concludes.
\todo{Update once finished}